% !TEX root = ../main.tex

\section{Introduction}
\label{sec:introduction}

Human visual attention is not distributed uniformly over a scene: observers consistently fixate on a small, predictable subset of regions. \textit{Visual saliency prediction} estimates, for a given image, the distribution of human eye fixations as a continuous \textit{saliency map}. Reliable saliency models support perceptually-driven image/video compression, thumbnailing, retargeting, advertising design, and attention priors for downstream vision and robotics pipelines.

Since the introduction of large fixation datasets such as MIT1003~\cite{judd2009learning} and SALICON~\cite{jiang2015salicon}, the field has moved from hand-crafted features to end-to-end CNN encoder-decoder models, typically initialized from ImageNet-pretrained backbones~\cite{wang2018deep,ghariba2019visual}; features learned for object recognition have repeatedly been shown to transfer well to saliency~\cite{kummerer2016deepgaze}. More recently, Transformer self-attention has been proposed to capture long-range dependencies that convolutional receptive fields struggle with~\cite{lou2022transalnet}. \secref{sec:related_work} discusses these works in more detail.

This report addresses two related problems. The \textbf{technical} one: designing, training and evaluating a saliency network under a tight compute budget (a single Colab T4 GPU, \mytexttilde15\,GB VRAM), starting from a from-scratch convolutional encoder-decoder and evolving it into a hybrid model with a pretrained Transformer encoder, while keeping the decoder and loss fixed for a controlled comparison. The \textbf{process} one: how a small team with no prior hands-on experience in saliency prediction iterated from a naive replica to a cross-dataset validated model. We keep this development trajectory visible (\secref{sec:workflow}) because the reasoning behind each design decision is, in our experience, as instructive as the final architecture.

\noindent \textbf{Contributions.} This report: (i) describes a multi-decoder, deep-supervised saliency network trained from scratch on SALICON with a combined KL-divergence/correlation/similarity loss applied to the fused output and three intermediate decoder branches (\secref{sec:model11}); (ii) extends it by replacing the encoder with a pretrained Swin-Tiny Transformer while keeping decoder, fusion and loss fixed, isolating the effect of the encoder swap (\secref{sec:model20}); (iii) evaluates both cross-dataset on MIT1003, discussing where the expected Transformer gains did, and did not, materialize (\secref{sec:validation}, \secref{sec:testing}); (iv) documents the iterative engineering process behind these models (\secref{sec:workflow}).

\noindent \textbf{Structure.} \secref{sec:related_work} positions our work in the literature. \secref{sec:workflow} describes the development process. \secref{sec:model11} and \secref{sec:model20} detail the CNN and Transformer-encoder models, respectively. \secref{sec:validation} and \secref{sec:testing} describe the evaluation protocol and preliminary results, and \secref{sec:conclusions} discusses findings and open work.
